% ============================================================
%  MTH636 – Game Theory | Lecture 32
%  Agent Form Representation, Correlated Equilibrium &
%  Introduction to Auction Theory
% ============================================================
\documentclass[12pt, a4paper]{article}

% ============================================================
%  MTH636 -- Game Theory
%  Notes Preamble — Color-Coded, TikZ-Enhanced
% ============================================================

% --- Core packages ---
\usepackage[T1]{fontenc}
\usepackage[utf8]{inputenc}
\usepackage{lmodern}
\usepackage[margin=2.5cm, top=2.8cm, bottom=2.8cm]{geometry}
\usepackage{microtype}

% --- Math ---
\usepackage{amsmath, amssymb, amsthm, mathtools}
\usepackage{bm}        % bold math
\usepackage{cancel}    % strike-through in math

% --- Color & Boxes ---
\usepackage[dvipsnames, svgnames, x11names, table]{xcolor}
\usepackage[most]{tcolorbox}
\tcbuselibrary{breakable, skins, theorems}

% --- TikZ & Diagrams ---
\usepackage{tikz}
\usepackage{tikz-qtree}
\usetikzlibrary{arrows.meta, positioning, shapes.geometric, fit, backgrounds,
                decorations.pathreplacing, calc, matrix, trees}
\usepackage{pgfplots}
\pgfplotsset{compat=1.18}

% --- Tables ---
\usepackage{booktabs}
\usepackage{multirow}
\usepackage{array}
\usepackage{tabularx}
\usepackage{makecell}

% --- Typography & Layout ---
\usepackage{enumitem}
\usepackage{parskip}
\usepackage{fancyhdr}
\usepackage{titlesec}
\usepackage{titletoc}
\usepackage{hyperref}
\usepackage{graphicx}
\usepackage{float}
\usepackage{caption}

% --- Color Palette ---
\definecolor{defblue}{HTML}{1A5276}      % Definition header
\definecolor{defbluebg}{HTML}{D6EAF8}    % Definition background
\definecolor{thmgreen}{HTML}{1E8449}     % Theorem header
\definecolor{thmbg}{HTML}{D5F5E3}        % Theorem background
\definecolor{exorange}{HTML}{B7770D}     % Example header
\definecolor{exbg}{HTML}{FEF9E7}         % Example background
\definecolor{warnred}{HTML}{922B21}      % Warning/Key Insight header
\definecolor{warnbg}{HTML}{FADBD8}       % Warning background
\definecolor{noticegray}{HTML}{2C3E50}   % Note header
\definecolor{noticebg}{HTML}{EBF5FB}     % Note background
\definecolor{taketawayyellow}{HTML}{7D6608} % Key Takeaway header
\definecolor{takeawaybg}{HTML}{FDFEFE}   % Takeaway background
\definecolor{accentpurple}{HTML}{6C3483} % Accent / proof
\definecolor{lightgray}{HTML}{F2F3F4}    % shading

% ============================================================
%  TCOLORBOX ENVIRONMENTS
% ============================================================

% -- Definition (Blue) --
\newtcolorbox{defbox}[2][]{
  enhanced, breakable,
  colback=defbluebg, colframe=defblue, coltitle=white,
  fonttitle=\bfseries\sffamily,
  title={Definition: #2},
  attach boxed title to top left={yshift=-2mm, xshift=4mm},
  boxed title style={colback=defblue, rounded corners},
  arc=4pt, boxrule=1.2pt,
  left=6pt, right=6pt, top=8pt, bottom=6pt,
  #1
}

% -- Theorem / Proposition (Green) --
\newtcolorbox{thmbox}[2][]{
  enhanced, breakable,
  colback=thmbg, colframe=thmgreen, coltitle=white,
  fonttitle=\bfseries\sffamily,
  title={Theorem: #2},
  attach boxed title to top left={yshift=-2mm, xshift=4mm},
  boxed title style={colback=thmgreen, rounded corners},
  arc=4pt, boxrule=1.2pt,
  left=6pt, right=6pt, top=8pt, bottom=6pt,
  #1
}

% -- Proposition (Green variant) --
\newtcolorbox{propbox}[2][]{
  enhanced, breakable,
  colback=thmbg, colframe=thmgreen!70!black, coltitle=white,
  fonttitle=\bfseries\sffamily,
  title={Proposition: #2},
  attach boxed title to top left={yshift=-2mm, xshift=4mm},
  boxed title style={colback=thmgreen!70!black, rounded corners},
  arc=4pt, boxrule=1.2pt,
  left=6pt, right=6pt, top=8pt, bottom=6pt,
  #1
}

% -- Example (Orange/Yellow) --
\newtcolorbox{exbox}[2][]{
  enhanced, breakable,
  colback=exbg, colframe=exorange, coltitle=white,
  fonttitle=\bfseries\sffamily,
  title={Example: #2},
  attach boxed title to top left={yshift=-2mm, xshift=4mm},
  boxed title style={colback=exorange, rounded corners},
  arc=4pt, boxrule=1.2pt,
  left=6pt, right=6pt, top=8pt, bottom=6pt,
  #1
}

% -- Key Insight / Warning (Red) --
\newtcolorbox{warnbox}[2][]{
  enhanced, breakable,
  colback=warnbg, colframe=warnred, coltitle=white,
  fonttitle=\bfseries\sffamily,
  title={Key Insight: #2},
  attach boxed title to top left={yshift=-2mm, xshift=4mm},
  boxed title style={colback=warnred, rounded corners},
  arc=4pt, boxrule=1.2pt,
  left=6pt, right=6pt, top=8pt, bottom=6pt,
  #1
}

% -- Note (Gray-Blue) --
\newtcolorbox{notebox}[1][]{
  enhanced, breakable,
  colback=noticebg, colframe=noticegray!60, coltitle=white,
  fonttitle=\bfseries\sffamily,
  title={Note},
  attach boxed title to top left={yshift=-2mm, xshift=4mm},
  boxed title style={colback=noticegray!60, rounded corners},
  arc=4pt, boxrule=0.8pt,
  left=6pt, right=6pt, top=8pt, bottom=6pt,
  #1
}

% -- Key Takeaway --
\newtcolorbox{takebox}[1][]{
  enhanced, breakable,
  colback=takeawaybg, colframe=taketawayyellow, coltitle=white,
  fonttitle=\bfseries\sffamily,
  title={Key Takeaway},
  attach boxed title to top left={yshift=-2mm, xshift=4mm},
  boxed title style={colback=taketawayyellow, rounded corners},
  arc=4pt, boxrule=1.2pt,
  left=6pt, right=6pt, top=8pt, bottom=6pt,
  #1
}

% -- Intuition --
\newtcolorbox{intubox}[1][]{
  enhanced, breakable,
  colback=accentpurple!8, colframe=accentpurple!60, coltitle=white,
  fonttitle=\bfseries\sffamily,
  title={Intuition},
  attach boxed title to top left={yshift=-2mm, xshift=4mm},
  boxed title style={colback=accentpurple!80, rounded corners},
  arc=4pt, boxrule=0.8pt,
  left=6pt, right=6pt, top=8pt, bottom=6pt,
  #1
}

% ============================================================
%  PAYOFF MATRIX HELPER
%  Usage: inside tabular, use \payoff{a}{b} for (a,b)
% ============================================================
\newcommand{\payoff}[2]{$(#1,\;#2)$}
\newcommand{\nash}[2]{\underline{$#1$},\underline{$#2$}}  % underline NE payoffs
\newcommand{\Nset}{\mathbb{N}}
\newcommand{\Rset}{\mathbb{R}}
\newcommand{\Sset}{\mathcal{S}}
\newcommand{\Ui}{u_i}

% ============================================================
%  SECTION FORMATTING
% ============================================================
\titleformat{\section}
  {\large\bfseries\sffamily\color{defblue}}
  {\thesection}{1em}{}[\titlerule]

\titleformat{\subsection}
  {\normalsize\bfseries\sffamily\color{thmgreen!80!black}}
  {\thesubsection}{1em}{}

\titleformat{\subsubsection}
  {\normalsize\itshape\sffamily\color{exorange!80!black}}
  {\thesubsubsection}{1em}{}

% ============================================================
%  HEADER / FOOTER
% ============================================================
\pagestyle{fancy}
\fancyhf{}
\fancyhead[L]{\sffamily\small\color{defblue} MTH636 -- Game Theory}
\fancyhead[R]{\sffamily\small\color{defblue} \leftmark}
\fancyfoot[C]{\sffamily\small\thepage}
\renewcommand{\headrulewidth}{0.4pt}

% ============================================================
%  HYPERREF SETUP
% ============================================================
\hypersetup{
  colorlinks=true,
  linkcolor=defblue,
  urlcolor=thmgreen,
  citecolor=accentpurple,
  bookmarksnumbered=true,
  pdftitle={MTH636 Game Theory Notes},
  pdfauthor={MTH636 Student Notes}
}

% ============================================================
%  MISC HELPERS
% ============================================================
\newcommand{\BR}{\text{BR}}        % Best Response
\newcommand{\NE}{\text{NE}}        % Nash Equilibrium
\newcommand{\SPNE}{\text{SPNE}}    % Subgame Perfect NE
\newcommand{\Dom}{\text{dom}}      % Dominance
\newcommand{\mixs}[1]{\sigma_{#1}} % Mixed strategy
\newcommand{\pures}[1]{s_{#1}}     % Pure strategy


% ============================================================
\begin{document}

% Title
\begin{center}
  {\LARGE\bfseries\sffamily\color{defblue} MTH636 -- Game Theory}\\[6pt]
  {\Large\sffamily Lecture 32: Agent Form Games, Correlated Equilibrium \& Auction Theory}\\[4pt]
\end{center}
\vspace{0.5cm}
\hrule height 1.5pt
\vspace{0.3cm}

% ============================================================
\section{Agent Form Representation \& Bayesian Nash Equilibrium}

In Lecture 31, we introduced Harsanyi games of incomplete information and defined the concept of \textbf{Bayesian Nash Equilibrium (BNE)}. We now establish the formal mathematical connection between a Harsanyi game and a strategic form game called the \textbf{Agent Form Game}, which will allow us to prove both the existence of BNE and its fundamental equivalence with Nash Equilibrium.

\subsection{Formulation of the Agent Form Game}

Let $\Gamma = \langle N, (T_i)_{i \in N}, p, S, (u_t)_{t \in T} \rangle$ be an incomplete information game, where:
\begin{itemize}
    \item $N = \{1, 2, \dots, n\}$ is the set of players.
    \item $T_i$ is the finite type set of player $i$, and $T = \prod_{i=1}^n T_i$.
    \item $p$ is the common prior probability distribution over $T$ with $p(t) > 0$ for all $t \in T$.
    \item $A_i(t_i)$ is the action set available to player $i$ when their type is $t_i$.
    \item $u_i(t, a)$ is player $i$'s payoff given type profile $t \in T$ and action profile $a \in \prod_i A_i(t_i)$.
\end{itemize}

We formulate an associated strategic form game called the \textbf{Agent Form Game}, denoted by $\hat{\Gamma}$.

\begin{defbox}{Agent Form Game $\hat{\Gamma}$}
The \textbf{Agent Form Game} associated with the incomplete information game $\Gamma$ is defined as the tuple:
\[
  \hat{\Gamma} = \langle \hat{N}, (\hat{S}_\ell)_{\ell \in \hat{N}}, (\hat{u}_\ell)_{\ell \in \hat{N}} \rangle
\]
where:
\begin{enumerate}
    \item \textbf{Set of Agents:} Each type of each player is treated as an autonomous decision-maker:
    \[
      \hat{N} = \bigcup_{i=1}^n T_i
    \]
    Thus, an agent $\ell \in \hat{N}$ corresponds to a specific pair $(i, t_i)$ where $t_i \in T_i$.
    \item \textbf{Strategy Sets:} If agent $\ell$ corresponds to type $t_i \in T_i$ of player $i$, their pure strategy set is simply the action set of that type:
    \[
      \hat{S}_\ell = A_i(t_i) \quad \text{where } \ell = t_i
    \]
    \item \textbf{Payoff Functions:} When an action profile $a = (a_\ell)_{\ell \in \hat{N}}$ is chosen by all agents, the payoff of agent $\ell = t_i \in T_i$ is the \textbf{interim expected utility} of player $i$ conditional on having type $t_i$:
    \[
      \hat{u}_\ell(a) = \sum_{t_{-i} \in T_{-i}} p(t_{-i} \mid t_i) \, u_i\!\left( (t_i, t_{-i}),\, (a_{t_i}, (a_{t_j})_{j \neq i}) \right)
    \]
    where $p(t_{-i} \mid t_i) = \frac{p(t_i, t_{-i})}{\sum_{t'_{-i}} p(t_i, t'_{-i})}$ is the conditional probability of opponent types.
\end{enumerate}
\end{defbox}

\begin{notebox}
\textbf{Remark on Strategy Correspondences:}
For any player $i \in N$, a behavior strategy $\sigma_i$ specifies an action distribution $\sigma_i(t_i) \in \Delta(A_i(t_i))$ for each type $t_i \in T_i$. In the agent form game $\hat{\Gamma}$, each such distribution is precisely a mixed strategy for agent $\ell = t_i$. 

Conversely, any profile of mixed strategies for the agents $\{t_i\}_{t_i \in T_i}$ in the agent form game directly defines a valid behavior strategy for player $i$ in the incomplete information game $\Gamma$.
\end{notebox}

\subsection{Equivalence and Existence Theorems}

\begin{thmbox}{Equivalence of BNE and Agent Form NE}
A behavior strategy profile $(\sigma_i)_{i \in N}$ is a \textbf{Bayesian Nash Equilibrium (BNE)} in the incomplete information game $\Gamma$ if and only if the corresponding strategy vector $(\sigma_i(t_i))_{i \in N,\, t_i \in T_i}$ is a \textbf{Nash Equilibrium} in the agent form game $\hat{\Gamma}$.
\end{thmbox}

\begin{proof}
$(\implies)$ Suppose $(\sigma_i)_{i \in N}$ is a BNE of $\Gamma$. By definition, for every player $i \in N$ and every type $t_i \in T_i$:
\[
  U_i(\sigma \mid t_i) \geq U_i\!\left( (a_i, \sigma_{-i}) \mid t_i \right) \quad \forall\, a_i \in A_i(t_i)
\]
Since the payoff function of agent $\ell = t_i$ in $\hat{\Gamma}$ is defined identically as the interim expected utility $\hat{u}_\ell = U_i(\cdot \mid t_i)$, no agent $\ell$ can profitably deviate from $\sigma_i(t_i)$. Thus, the strategy vector is a Nash Equilibrium of $\hat{\Gamma}$.

$(\impliedby)$ We prove the contrapositive. Suppose $(\sigma_i)_{i \in N}$ is \textbf{not} a BNE of $\Gamma$. Then there exists some player $i \in N$, some type $t_i \in T_i$, and an action $a \in A_i(t_i)$ such that deviating to $a$ yields a strictly higher interim utility:
\[
  U_i(\sigma \mid t_i) < U_i\!\left( (a, \sigma_{-i}) \mid t_i \right)
\]
In the agent form game $\hat{\Gamma}$, this implies that agent $\ell = t_i$ can unilaterally deviate to action $a$ and achieve:
\[
  \hat{u}_\ell\!\left( (\sigma_j(t_j))_{j \in N,\, t_j \in T_j} \right) < \hat{u}_\ell\!\left( a,\, \sigma_{-\ell} \right)
\]
Hence, the strategy profile is not a Nash Equilibrium of $\hat{\Gamma}$.
\end{proof}

\begin{thmbox}{Existence of Bayesian Nash Equilibrium}
Every finite incomplete information game $\Gamma$ possesses at least one \textbf{Bayesian Nash Equilibrium} in behavior strategies.
\end{thmbox}

\begin{proof}
The agent form game $\hat{\Gamma}$ is a standard finite strategic form game with player set $\hat{N}$ and pure action sets $\hat{S}_\ell$. By Nash's Existence Theorem (1950), every finite normal form game has at least one Nash Equilibrium in mixed strategies. By the Equivalence Theorem above, this mixed Nash Equilibrium in $\hat{\Gamma}$ is precisely a Bayesian Nash Equilibrium in behavior strategies for $\Gamma$.
\end{proof}

\begin{thmbox}{BNE and Strategic Form Nash Equilibrium}
In any Harsanyi incomplete information model with strictly positive common prior $p(t) > 0$ for all $t \in T$, a strategy profile is an ex-ante Nash Equilibrium if and only if it is a Bayesian Nash Equilibrium.
\end{thmbox}

\subsection{Worked Example: Constructing the Agent Form Game}

\begin{exbox}{Agent Form Representation with Asymmetric Types}
Consider an incomplete information game with two players, $N = \{1, 2\}$:
\begin{itemize}
    \item Player 1 has two types: $T_1 = \{t_1, t_2\}$ (often denoted $l_1, l_2$).
    \item Player 2 has one type: $T_2 = \{\text{II}\}$.
    \item Common prior probabilities: $p(t_1, \text{II}) = \frac{1}{2}$ and $p(t_2, \text{II}) = \frac{1}{2}$.
\end{itemize}
Conditional beliefs are trivially:
\[
  p(\text{II} \mid t_1) = 1, \quad p(\text{II} \mid t_2) = 1, \quad p(t_1 \mid \text{II}) = \frac{1/2}{1} = \frac{1}{2}, \quad p(t_2 \mid \text{II}) = \frac{1}{2}
\]

The state payoff matrices are:
\begin{center}
\begin{tabular}{cc}
\textbf{State $(t_1, \text{II})$} & \textbf{State $(t_2, \text{II})$} \\[4pt]
\begin{tabular}{c|cc}
 & $L$ & $R$ \\ \hline
$T_1$ & $(1, 0)$ & $(0, 2)$ \\
$B_1$ & $(0, 3)$ & $(1, 0)$
\end{tabular}
&
\begin{tabular}{c|cc}
 & $L$ & $R$ \\ \hline
$T_2$ & $(0, 2)$ & $(1, 1)$ \\
$B_2$ & $(1, 0)$ & $(0, 2)$
\end{tabular}
\end{tabular}
\end{center}

\textbf{Step 1: Formulate the Agent Form Game $\hat{\Gamma}$:}
\begin{itemize}
    \item Set of agents: $\hat{N} = \{t_1, t_2, \text{II}\}$ (3 agents).
    \item Action sets: $\hat{S}_{t_1} = \{T_1, B_1\}$, $\hat{S}_{t_2} = \{T_2, B_2\}$, $\hat{S}_{\text{II}} = \{L, R\}$.
\end{itemize}

\textbf{Step 2: Payoff Computations for Profiles:}
Consider the pure action profile $(T_1, B_2, L)$:
\begin{itemize}
    \item For agent $t_1$:
    \[
      \hat{u}_{t_1}(T_1, B_2, L) = 1 \cdot u_1((t_1, \text{II}), (T_1, L)) = 1
    \]
    \item For agent $t_2$:
    \[
      \hat{u}_{t_2}(T_1, B_2, L) = 1 \cdot u_1((t_2, \text{II}), (B_2, L)) = 1
    \]
    \item For agent $\text{II}$:
    \begin{align*}
      \hat{u}_{\text{II}}(T_1, B_2, L) &= u_2((t_1, \text{II}), (T_1, L)) \, p(t_1 \mid \text{II}) + u_2((t_2, \text{II}), (B_2, L)) \, p(t_2 \mid \text{II}) \\
      &= 0 \cdot \frac{1}{2} + 0 \cdot \frac{1}{2} = 0
    \end{align*}
\end{itemize}

\textbf{Step 3: Verification of Deviation Equivalence:}
Let mixed strategies be:
\[
  \sigma_{t_1} = x [T_1] + (1-x)[B_1], \quad \sigma_{t_2} = y [T_2] + (1-y)[B_2], \quad \sigma_{\text{II}} = q [L] + (1-q)[R]
\]
Suppose this strategy profile is not a BNE because Player 1 can profitably deviate at type $t_1$ by playing pure action $T_1$. Then:
\[
  U_1(\sigma \mid t_1) < U_1((T_1, \sigma_{\text{II}}) \mid t_1)
\]
The ex-ante utility of Player 1 decomposes as:
\[
  U_1(\sigma) = p(t_1) \, U_1(\sigma \mid t_1) + p(t_2) \, U_1(\sigma \mid t_2)
\]
Defining the modified strategy $\bar{\sigma}_1$ by:
\[
  \bar{\sigma}_1(t) = \begin{cases} T_1 & \text{if } t = t_1 \\ y[T_2] + (1-y)[B_2] & \text{if } t = t_2 \end{cases}
\]
we immediately obtain $U_1(\sigma) < U_1(\bar{\sigma}_1, \sigma_{\text{II}})$, establishing that interim optimality per agent is strictly equivalent to ex-ante optimality.
\end{exbox}

% ============================================================
\newpage
\section{Correlated Equilibrium}

The concept of \textbf{Correlated Equilibrium}, introduced by Robert Aumann in 1974, generalizes Nash Equilibrium by allowing players to coordinate their choices through public or private signals generated by a correlation device (such as a mediator or randomizing mechanism).

\subsection{Motivating Example: The Traffic Intersection (Chicken Game)}

Consider two cars approaching an intersection simultaneously. Each driver can either \textbf{Cross ($C$)} or \textbf{Stop ($S$)}.

\begin{center}
\renewcommand{\arraystretch}{1.2}
\begin{tabular}{c|cc}
 & $C$ & $S$ \\ \hline
$C$ & $(-10, -10)$ & $(5, 0)$ \\
$S$ & $(0, 5)$ & $(-1, -1)$
\end{tabular}
\end{center}

\subsubsection{Standard Nash Equilibrium Analysis}
\begin{itemize}
    \item \textbf{Pure Strategy NE:} There are two asymmetric pure Nash equilibria:
    \[
      (C, S) \implies \text{payoffs } (5, 0), \qquad (S, C) \implies \text{payoffs } (0, 5)
    \]
    Both are inherently unfair to one driver, and in a simultaneous game with no prior communication, players cannot determine which equilibrium to coordinate upon.
    \item \textbf{Symmetric Mixed Strategy NE:} Let Player 2 choose Cross with probability $q$ and Stop with probability $1-q$. For Player 1 to be indifferent:
    \[
      U_1(C) = U_1(S) \implies -10q + 5(1-q) = 0q - 1(1-q)
    \]
    \[
      5 - 15q = q - 1 \implies 16q = 6 \implies q = \frac{3}{8}
    \]
    Thus, the unique symmetric mixed strategy NE is:
    \[
      \sigma_1^* = \sigma_2^* = \left( \frac{3}{8}[C] + \frac{5}{8}[S] \right)
    \]
    Under this mixed equilibrium, the expected utility of each player is:
    \[
      U_i(\sigma^*) = \frac{3}{8}(0) + \frac{5}{8}(-1) = -\frac{5}{8} = -0.625
    \]
    Both players suffer negative expected payoffs, and a disastrous crash $(C, C)$ occurs with probability $\left(\frac{3}{8}\right)^2 = \frac{9}{64} \approx 14.1\%$.
\end{itemize}

\subsubsection{Introducing a Correlation Device (Traffic Light / Mediator)}

Now suppose an external device (e.g., a traffic signal) selects a recommended strategy profile $(s_1, s_2)$ according to the following joint probability distribution $p$:

\begin{center}
\renewcommand{\arraystretch}{1.2}
\begin{tabular}{c|cc|c}
 & $C$ & $S$ & Total \\ \hline
$C$ & $0$ & $0.55$ & $0.55$ \\
$S$ & $0.40$ & $0.05$ & $0.45$ \\ \hline
Total & $0.40$ & $0.60$ & $1.00$
\end{tabular}
\end{center}

The recommendations are \textbf{private}: Player 1 is told only their own action $s_1$, and Player 2 is told only their own action $s_2$. Both players know the joint distribution $p$.

\textbf{Checking Player 1's Incentive to Obey:}
\begin{enumerate}
    \item \textbf{Case 1: Player 1 is recommended Cross ($s_1 = C$):}
    \begin{itemize}
        \item Since $p(C, C) = 0$, Player 1 updates their belief about Player 2:
        \[
          p(s_2 = S \mid s_1 = C) = \frac{0.55}{0 + 0.55} = 1
        \]
        Player 1 knows with certainty that Player 2 was told to Stop ($S$).
        \item If Player 1 obeys and plays $C$: Payoff is $u_1(C, S) = 5$.
        \item If Player 1 deviates to Stop ($S$): Payoff is $u_1(S, S) = -1$.
        \item Since $5 > -1$, Player 1 has \textbf{no incentive to deviate}.
    \end{itemize}
    
    \item \textbf{Case 2: Player 1 is recommended Stop ($s_1 = S$):}
    \begin{itemize}
        \item The marginal probability of getting recommendation $S$ is $p(s_1 = S) = 0.40 + 0.05 = 0.45$.
        \item Conditional probability that Player 2 received $C$:
        \[
          p(s_2 = C \mid s_1 = S) = \frac{0.40}{0.45} = \frac{8}{9} \approx 0.89
        \]
        \item Conditional probability that Player 2 received $S$:
        \[
          p(s_2 = S \mid s_1 = S) = \frac{0.05}{0.45} = \frac{1}{9} \approx 0.11
        \]
        \item If Player 1 obeys the recommendation and plays $S$:
        \[
          \mathbb{E}[u_1 \mid S] = \frac{8}{9}(0) + \frac{1}{9}(-1) = -0.11
        \]
        \item If Player 1 disobeys and plays $C$:
        \[
          \mathbb{E}[u_1 \mid C] = \frac{8}{9}(-10) + \frac{1}{9}(5) = -\frac{80}{9} + \frac{5}{9} = -\frac{75}{9} \approx -8.33
        \]
        \item Since $-0.11 > -8.33$, Player 1 strictly prefers to \textbf{obey the signal and stop}.
    \end{itemize}
\end{enumerate}

\begin{takebox}
\textbf{Payoff Comparison:}
Under this correlated distribution $p$, Player 1's expected payoff is:
\[
  U_1(p) = 0.55(5) + 0.40(0) + 0.05(-1) = 2.75 - 0.05 = 2.70
\]
Notice that $2.70 \gg -\frac{5}{8} = -0.625$. Both players achieve substantially higher expected payoffs than the mixed NE, collisions are completely avoided ($p(C, C) = 0$), and yet neither player has any incentive to disobey the mediator!
\end{takebox}

\subsection{Formal Definition of Correlated Equilibrium}

Let $G = \langle N, (S_i)_{i \in N}, (u_i)_{i \in N} \rangle$ be a strategic form game, and let $S = \prod_{i=1}^n S_i$.

\begin{defbox}{Extensive Form Game $\Gamma^*(p)$ of a Correlation Device}
Let $p \in \Delta(S)$ be a probability distribution over the pure strategy space $S$. The correlation device defines an extensive form game $\Gamma^*(p)$ as follows:
\begin{enumerate}
    \item Nature chooses a strategy profile $s = (s_1, \dots, s_n) \in S$ according to distribution $p$.
    \item Each player $i \in N$ is privately informed of their recommended action $s_i$ (and learns nothing about $s_{-i}$).
    \item Each player $i$ then chooses an action $s'_i \in S_i$.
    \item Payoffs are given by $u_i(s'_1, \dots, s'_n)$.
\end{enumerate}
In $\Gamma^*(p)$, a pure strategy for player $i$ is a mapping $\tau_i : S_i \to S_i$, where $\tau_i(s_i)$ is the action taken when recommended $s_i$. The strategy of \textbf{following recommendations} is denoted by:
\[
  \tilde{\tau}_i(s_i) = s_i \quad \forall\, s_i \in S_i
\]
\end{defbox}

\begin{defbox}{Correlated Equilibrium (Aumann, 1974)}
A probability distribution $p \in \Delta(S)$ is a \textbf{Correlated Equilibrium} of the game $G$ if the obedience strategy profile $\tilde{\tau} = (\tilde{\tau}_i)_{i \in N}$ is a Nash Equilibrium of the game $\Gamma^*(p)$.
\end{defbox}

\begin{warnbox}{Privacy of Recommendations}
It is critical that recommendations are \textbf{private}. If Player 1 knew Player 2's recommendation in advance, the strategic interaction would collapse back into standard subgames, potentially destroying the incentive to obey.
\end{warnbox}

\subsection{Characterization Theorem and Proof}

\begin{thmbox}{Linear Incentive Constraints for Correlated Equilibrium}
A probability distribution $p \in \Delta(S)$ is a Correlated Equilibrium of $G$ if and only if for every player $i \in N$ and for all actions $s_i, s'_i \in S_i$:
\[
  \sum_{s_{-i} \in S_{-i}} p(s_i, s_{-i}) \, u_i(s_i, s_{-i}) \;\geq\; \sum_{s_{-i} \in S_{-i}} p(s_i, s_{-i}) \, u_i(s'_i, s_{-i})
\]
\end{thmbox}

\begin{proof}
Fix player $i \in N$, a recommended strategy $s_i \in S_i$, and an alternative strategy $s'_i \in S_i$.

\textbf{Case 1:} Suppose $\sum_{s_{-i} \in S_{-i}} p(s_i, s_{-i}) = 0$. Then $p(s_i, s_{-i}) = 0$ for all $s_{-i}$. In this case, recommendation $s_i$ occurs with probability zero, both sides of the inequality evaluate to 0, and the condition is trivially satisfied.

\textbf{Case 2:} Suppose $\sum_{s_{-i} \in S_{-i}} p(s_i, s_{-i}) > 0$. By Bayes' Rule, when player $i$ receives recommendation $s_i$, their posterior belief about opponent actions is:
\[
  p(s_{-i} \mid s_i) = \frac{p(s_i, s_{-i})}{\sum_{t_{-i} \in S_{-i}} p(s_i, t_{-i})}
\]
The expected utility of player $i$ from following the recommendation $s_i$ (assuming all opponents also follow their recommendations) is:
\[
  \mathbb{E}[u_i \mid s_i] = \sum_{s_{-i} \in S_{-i}} p(s_{-i} \mid s_i) \, u_i(s_i, s_{-i})
\]
If player $i$ instead deviates to action $s'_i$, their expected utility is:
\[
  \mathbb{E}[u_i \mid s'_i] = \sum_{s_{-i} \in S_{-i}} p(s_{-i} \mid s_i) \, u_i(s'_i, s_{-i})
\]
By definition, $\tilde{\tau}$ is a Nash Equilibrium of $\Gamma^*(p)$ iff no player has an incentive to deviate upon receiving any recommendation $s_i$:
\[
  \sum_{s_{-i} \in S_{-i}} p(s_{-i} \mid s_i) \, u_i(s_i, s_{-i}) \;\geq\; \sum_{s_{-i} \in S_{-i}} p(s_{-i} \mid s_i) \, u_i(s'_i, s_{-i}) \quad \forall\, i \in N,\ \forall\, s_i, s'_i \in S_i
\]
Substituting the formula for $p(s_{-i} \mid s_i)$:
\[
  \frac{\sum_{s_{-i}} p(s_i, s_{-i}) \, u_i(s_i, s_{-i})}{\sum_{t_{-i}} p(s_i, t_{-i})} \;\geq\; \frac{\sum_{s_{-i}} p(s_i, s_{-i}) \, u_i(s'_i, s_{-i})}{\sum_{t_{-i}} p(s_i, t_{-i})}
\]
Multiplying both sides by the strictly positive denominator $\sum_{t_{-i}} p(s_i, t_{-i}) = p(s_i)$ yields:
\[
  \sum_{s_{-i} \in S_{-i}} p(s_i, s_{-i}) \, u_i(s_i, s_{-i}) \;\geq\; \sum_{s_{-i} \in S_{-i}} p(s_i, s_{-i}) \, u_i(s'_i, s_{-i})
\]
which is the desired characterization.
\end{proof}

\subsection{Properties of Correlated Equilibria}

\begin{thmbox}{Mixed NE Induces a Correlated Equilibrium}
Every Nash Equilibrium (pure or mixed) of a strategic form game $G$ induces a Correlated Equilibrium.
\end{thmbox}

\begin{proof}
Let $\sigma^* = (\sigma_1^*, \dots, \sigma_n^*)$ be a Nash Equilibrium of $G$. Define the probability distribution $p_{\sigma^*}$ on $S$ as the product distribution:
\[
  p_{\sigma^*}(s_1, \dots, s_n) = \sigma_1^*(s_1) \cdot \sigma_2^*(s_2) \cdots \sigma_n^*(s_n)
\]
Since the distribution factors independently, knowing $s_i$ provides no information about $s_{-i}$:
\[
  p(s_{-i} \mid s_i) = \sigma_{-i}^*(s_{-i})
\]
The obedience condition for player $i$ becomes:
\[
  \sum_{s_{-i}} \sigma_{-i}^*(s_{-i}) \, u_i(s_i, s_{-i}) \;\geq\; \sum_{s_{-i}} \sigma_{-i}^*(s_{-i}) \, u_i(s'_i, s_{-i})
\]
which is precisely the definition that $s_i$ is a best response to $\sigma_{-i}^*$ for every $s_i$ in the support of $\sigma_i^*$. Thus, $p_{\sigma^*}$ is a Correlated Equilibrium.
\end{proof}

\begin{intubox}
\textbf{Geometric Insight:}
The set of Correlated Equilibria is defined by a finite system of linear inequalities together with the probability simplex constraints ($\sum_s p(s) = 1$, $p(s) \geq 0$). 
Therefore, \textbf{the set of Correlated Equilibria is a convex polytope}. 
In contrast, the set of Nash Equilibria is generally non-convex and disconnected due to the non-linear independence constraints $p(s) = \prod_i \sigma_i(s_i)$.
\end{intubox}

\begin{exbox}{Correlated Equilibrium in the Prisoner's Dilemma}
Consider the standard Prisoner's Dilemma:
\begin{center}
\renewcommand{\arraystretch}{1.2}
\begin{tabular}{c|cc}
 & $C$ & $D$ \\ \hline
$C$ & $(3, 3)$ & $(0, 5)$ \\
$D$ & $(5, 0)$ & $(1, 1)$
\end{tabular}
\end{center}
\textbf{Question:} Is there any Correlated Equilibrium where the cooperative outcome $(C, C)$ receives positive probability?

\textbf{Analysis:} Suppose $p(C, C) > 0$. Then Player 1 can be recommended $C$ with positive probability. If Player 1 is recommended $C$, her expected payoff from obeying is:
\[
  \mathbb{E}[u_1 \mid C] = 3 \cdot p(C \mid C) + 0 \cdot p(D \mid C) = 3 \, p(C \mid C)
\]
If Player 1 deviates to $D$:
\[
  \mathbb{E}[u_1 \mid D] = 5 \cdot p(C \mid C) + 1 \cdot p(D \mid C) = 5 \, p(C \mid C) + p(D \mid C)
\]
Subtracting the two:
\[
  \mathbb{E}[u_1 \mid D] - \mathbb{E}[u_1 \mid C] = 2 \, p(C \mid C) + p(D \mid C) > 0
\]
Because $D$ strictly dominates $C$, Player 1 has a strictly positive incentive to defect regardless of any posterior beliefs about Player 2! 

Consequently, \textbf{no correlated equilibrium can assign positive probability to any strictly dominated strategy}. The unique correlated equilibrium of the Prisoner's Dilemma is $p(D, D) = 1$.
\end{exbox}

% ============================================================
\newpage
\section{Introduction to Auction Theory}

Auction theory models markets where goods are allocated via bidding mechanisms under asymmetric and incomplete information.

\subsection{Mathematical Preliminaries: Random Variables}

\begin{defbox}{Cumulative Distribution Function (CDF)}
For a real-valued random variable $X$, its \textbf{Cumulative Distribution Function (CDF)}, denoted $F_X : \mathbb{R} \to [0, 1]$, is defined by:
\[
  F_X(c) = \mathbb{P}(X \leq c) \quad \forall\, c \in \mathbb{R}
\]
A CDF is non-decreasing, right-continuous, and satisfies $\lim_{c \to -\infty} F_X(c) = 0$ and $\lim_{c \to \infty} F_X(c) = 1$.
\end{defbox}

\begin{defbox}{Probability Density Function (PDF)}
If $X$ is an absolutely continuous random variable, there exists a non-negative function $f_X : \mathbb{R} \to \mathbb{R}^+$, called the \textbf{Probability Density Function (PDF)}, such that:
\[
  F_X(c) = \int_{-\infty}^c f_X(x) \, dx \quad \forall\, c \in \mathbb{R}
\]
\end{defbox}

\subsection{Sealed-Bid Auctions with Independent Private Valuations (IPV)}

In the \textbf{Independent Private Valuations (IPV)} model, each buyer knows their own willingness to pay for the good, but does not know the valuations of the competing bidders.

\begin{defbox}{Sealed-Bid Auction Model $\langle N, (V_i, F_i)_{i \in N}, p, c \rangle$}
A sealed-bid auction with independent private valuations is defined by the tuple:
\[
  \langle N, (V_i, F_i)_{i \in N}, p, c \rangle
\]
where:
\begin{enumerate}
    \item $N = \{1, 2, \dots, n\}$ is the set of buyers (bidders).
    \item $V_i \subseteq \mathbb{R}$ is the set of all possible valuations of buyer $i$. Each buyer has a private valuation $v_i \in V_i$.
    \item $F_i$ is the Cumulative Distribution Function (CDF) from which buyer $i$'s valuation is drawn independently.
    \item $p : [0, \infty)^N \to \Delta N$ is the \textbf{allocation rule}, where $p_i(b)$ denotes the probability that buyer $i$ wins the object given the bid profile $b = (b_1, \dots, b_n)$. Here $\Delta N = \{ (p_1, \dots, p_n) \mid p_i \geq 0,\, \sum_{i=1}^n p_i \leq 1 \}$.
    \item $c : N \times [0, \infty)^N \to \mathbb{R}$ (or $c : [0, \infty)^N \to \mathbb{R}^N$) is the \textbf{payment rule}, determining the payment $c_i(b)$ that buyer $i$ must pay to the seller given the bid profile $b$.
\end{enumerate}
\end{defbox}

\subsection{Protocol of a Sealed-Bid Auction}

A sealed-bid auction proceeds through the following chronological stages:
\begin{enumerate}
    \item \textbf{Valuation Realization:} Nature draws each buyer's private valuation $v_i \in V_i$ independently according to distribution $F_i$.
    \item \textbf{Information Revelation:} Each buyer learns their own private valuation $v_i$, but does not observe the valuations $v_{-i}$ of any other buyers.
    \item \textbf{Bidding:} Every buyer simultaneously submits a bid $b_i \in [0, \infty)$ in a sealed/closed envelope.
    \item \textbf{Winner Determination:} The auctioneer opens all envelopes and selects the winner(s) according to the allocation rule $p(b)$. Typically, the highest bidder wins:
    \[
      p_i(b) = \begin{cases} 1 & \text{if } b_i > \max_{j \neq i} b_j \\ \frac{1}{k} & \text{if } b_i \text{ is tied among } k \text{ highest bidders} \\ 0 & \text{otherwise} \end{cases}
    \]
    \item \textbf{Payment Execution:} Buyers pay the seller according to payment rule $c(b)$. For example:
    \begin{itemize}
        \item \textbf{First-Price Auction:} The winning bidder pays their own bid, $c_i(b) = b_i$.
        \item \textbf{Second-Price (Vickrey) Auction:} The winning bidder pays the second-highest bid, $c_i(b) = \max_{j \neq i} b_j$.
    \end{itemize}
    \item \textbf{Utility Realization:} Assuming risk-neutral quasi-linear utility, buyer $i$'s realized utility is:
    \[
      u_i(v_i, b) = p_i(b) \cdot v_i - c_i(b)
    \]
\end{enumerate}

\end{document}
